% part2-defs.tex -- part 2 macros and styles (Proteek Rosul)
%   loaded by part2.tex and by ../combined.tex

% ======================================================================
%  3D SUPPORT
% ======================================================================
\usepackage{pgfplots}
\pgfplotsset{compat=1.18}

% overlay-aware visibility for tikz / pgfplots elements:
%   \addplot3[..., visible on=<2->] ...      \draw[visible on=<3->] ...
\tikzset{
  invisible/.style  = {opacity=0, text opacity=0},
  visible on/.style = {alt={#1{}{invisible}}},
  alt/.code args    = {<#1>#2#3}{\alt<#1>{\pgfkeysalso{#2}}{\pgfkeysalso{#3}}},
}

% height ramps: deep at the goal, pale on the rim.  Blue = Euclidean,
% orange = Manhattan -- the same pairing the 2D distance slide uses.
\pgfplotsset{
  colormap={hblue}{color=(rbluedk) color=(rbluemd) color=(rbluelt)},
  colormap={horange}{color=(cbOrange!70!black) color=(cbOrange) color=(cbOrangeLt)},
  % flat grid tiles, value = 0 free / 1 expanded / 2 start / 3 goal
  colormap={htiles}{color=(gfree) color=(gclose) color=(gstart) color=(ggoal)},
  land3d/.style   = {axis lines=none, clip=false},
  % land view={az}{el}{unit in cm}: pgfplots' own view={az}{el} projection,
  % but with one fixed length per data unit, so shapes keep their true
  % proportions (a cone stays round) and the size is predictable
  land view/.code n args = {3}{%
    \pgfmathsetmacro\lvxx{#3*cos(#1)}%
    \pgfmathsetmacro\lvxy{-#3*sin(#1)*sin(#2)}%
    \pgfmathsetmacro\lvyx{#3*sin(#1)}%
    \pgfmathsetmacro\lvyy{#3*cos(#1)*sin(#2)}%
    \pgfmathsetmacro\lvzy{#3*cos(#2)}%
    \edef\lvkeys{\noexpand\pgfkeysalso{%
      x={(\lvxx cm,\lvxy cm)}, y={(\lvyx cm,\lvyy cm)}, z={(0cm,\lvzy cm)}}}%
    \lvkeys},
  % lifted grid: flat-shaded facets with white gutters, like the 2D tiles
  tilesurf/.style = {surf, shader=faceted, faceted color=white,
                     line width=0.45pt, z buffer=sort, point meta=z},
}

\tikzset{
  sball/.style  = {circle, fill=gstart, draw=white, line width=0.7pt,
                   minimum size=3.4mm, inner sep=0pt},
  gball/.style  = {sball, fill=ggoal},
  vball/.style  = {sball, fill=hlyellow, draw=cbOrange, line width=1pt},
  tag/.style    = {font=\scriptsize, inner sep=1.5pt},
  tagbox/.style = {tag, fill=white, fill opacity=0.85, text opacity=1,
                   rounded corners=1pt},
  move3d/.style = {-{Stealth[length=2.2mm]}, line width=1.7pt,
                   line cap=round, shorten >=0.6pt},
}

% ---------- per-figure geometry ----------------------------------------
% 3D-a: 12 x 12 lattice, G at (7,3), height = kA * Euclidean distance
\def\kA{0.45}\def\gAx{7}\def\gAy{3}
\newcommand{\liftA}[2]{(#1,#2,{\kA*veclen((#1)-\gAx,(#2)-\gAy)})}
% the flat map on the floor (a loop body with # cannot sit inside a frame)
\newcommand{\floorA}{%
  \pgfplotsinvokeforeach{0,...,12}{%
    \draw[gridgray, line width=0.35pt] (##1,0,0) -- (##1,12,0);
    \draw[gridgray, line width=0.35pt] (0,##1,0) -- (12,##1,0);}}

% 3D-b: G at the origin, v at (6,0), height = kB * Euclidean distance
\def\kB{0.45}
\newcommand{\liftB}[2]{(#1,#2,{\kB*veclen(#1,#2)})}

% 3D-c: 18 x 12 cells, S and G cell centres, optimal cost C* = 5
%       Dijkstra expands the disk  |vS| <= C*;
%       A* the ellipse             |vS| + |vG| <= C* + slack
\def\kC{0.3}
\def\cSx{5.5}\def\cSy{6.5}\def\cGx{10.5}\def\cGy{6.5}
\def\cCost{5}\def\cSlack{0.9}
\newcommand{\legendC}{%
  \begin{tikzpicture}[x=0.34cm,y=0.34cm, baseline=0pt]
    \renewcommand{\gpad}{0.09}%
    \gcell{gclose}{0}{0}
      \node[anchor=west, font=\scriptsize] at (1.4,0.5) {expanded};
    \gcell{gfree}{7}{0}
      \node[anchor=west, font=\scriptsize] at (8.4,0.5) {never touched};
    \gcell{cbGreen}{16.5}{0}
      \node[anchor=west, font=\scriptsize] at (17.9,0.5) {optimal route};
  \end{tikzpicture}}

% 3D-e: cone / pyramid, rim at distance 4, sample point P = (2,1)
\def\kE{0.8}\def\pEx{2}\def\pEy{1}


% colour-coded highlight, shown from overlay 2 on: #1 = 2 orange, 3 green,
% 4 blue, 5 red, anything else yellow.  preamble2.tex calls it \hlon; it
% is renamed so it can sit beside part 1's \hlon{<overlay>}{...} in
% combined.tex.
\newcommand{\hlonc}[2]{%
  \alt<2->{%
    \setlength{\fboxsep}{1.4pt}%
    \ifnum#1=2 \colorbox{cbOrange!30}{#2}%
    \else\ifnum#1=3 \colorbox{cbGreen!30}{#2}%
    \else\ifnum#1=4 \colorbox{rblue!30}{#2}%
    \else\ifnum#1=5 \colorbox{cbRed!30}{#2}%
    \else \colorbox{hlyellow}{#2}%
    \fi\fi\fi\fi
  }{#2}%
}
